\documentstyle[% 


righttag% 


,11pt% 


,amssymb% 


,russian%               remove in Eng. 


,izvru%                 Set izven% in Eng. 


% ,izvru-ks             for brief communications 


]{amsart} 


 


%       Math definitions 


 


\newcommand{\la}{\langle} 


\newcommand{\ra}{\rangle} 


\newcommand{\Hom}{\operatorname{Hom}} 


\newcommand{\im}{\operatorname{Im}} 


\newcommand{\Ker}{\operatorname{Ker}} 


\newcommand{\tts}[1]{} 


 


%\hoffset=-15mm%                  For Eng. 


  


\setcounter{page}{42} 


\god{1998} 


\nomer{9~(436)} %                        Remove in Eng. 


%\nomer{Vol.\,42, No.\,9}%               Set in Eng. 


%\str{\pageref{begin}--\pageref{end}}%  Set in Eng. 


\udk{512.541} 


\author{\it С.Я.\,Гриншпон} 


% NB:   ^^ remove in Eng. 


\title{О равенстве нулю группы гомоморфизмов\\ 


абелевых групп} 


 


\date{} 


% \pagestyle{myheadings}%         Set in Eng. 


 


\pagestyle{plain} %              Remove in Eng. 


\begin{document} 


 


% \thispagestyle{plain}%          Set in Eng. 


 


\maketitle 


\label{begin} 


 


В работе решается вопрос об условиях, при которых группа гомоморфизмов 


абелевых групп $\Hom(A,C)$ равна нулю, если $C$ --- периодическая группа 


или однородная сепарабельная группа без кручения (в частности, группа 


без кручения ранга 1). Для периодической группы $C$ ответ на 


поставленный вопрос полностью получен. В случае групп без кручения $C$ 


получен полный ответ для групп $A$, у которых тип каждого ненулевого 


элемента их факторгрупп по периодической части не меньше типа группы $C$ 


(в частности, для однородных абелевых групп без кручения $A$, тип 


которых совпадает с типом группы $C$). Из полученных в данной работе 


результатов следует описание коузких абелевых групп, полученное в 


\cite{1}. 


 


Рассмотрим вначале общую ситуацию. Пусть $C$ --- произвольная абелева 


группа. Обозначим через $\frak A_C$ класс всех абелевых групп $A$ со 


свойством $\Hom(A,C)=0$. 


 


\begin{prop} 


Класс $\frak A_C$ замкнут относительно$:$ {\em а)}~факторгрупп$;$ {\em 


б)}~расширений$;$ {\em в)}~прямых сумм$;$ {\em г)}~прямых 


произведений с неизмеримым множеством компонент, если $C$ --- узкая 


группа$;$ {\em д)}~прямых пределов$;$ {\em е)}~тензорных произведений 


на произвольную абелеву группу. 


\end{prop} 


 


\begin{pf} 


а) Пусть $A\in\frak A_C$ и $B$ --- подгруппа группы $A$. Имеем 


$\Hom(A,C)=0$. Если предположить, 


что существует $\varphi\in\Hom(A/B,C)$, $\varphi\ne 0$, 


то, рассматривая 


гомоморфизм $\varphi\eta$, где $\eta$ --- 


канонический эпиморфизм $A$ на $A/B$, получим, что 


$\varphi\eta\in\Hom(A,C)$ и $\varphi\eta\ne 0$. Противоречие. 


 


б) Предположим, что в точной последовательности 


$0\to B\to A\to D\to 0$ группы $B$ и $D$ взяты 


из класса $\frak A_C$. Пусть $f:A\to C$ --- гомоморфизм. Отождествим 


группу $D$ с группой 


$A/B$. Построим гомоморфизм $\psi:A/B\to C$ следующим образом: 


$\psi(a+B)=f(a)$. Условие $\Hom(B,C)=0$ гарантирует корректность 


построенного отображения. Имеем $\varphi\eta=f$ ($\eta$ --- 


канонический 


эпиморфизм $A$ на $A/B$). Так как $A/B\in\frak A_C$, то $\psi=0$ и 


поэтому $f=0$. Значит, $A\in\frak A_C$. 


 


в) Пусть $\{A_i\}_{i\in I}$ --- некоторое семейство групп из класса 


$\frak A_C$. Тогда 


$\Hom\big(\underset{i\in I}{\oplus}{}A_i,C\big)\cong 


\prod\limits_{i\in I}\Hom(A_i,C)=0$. Следовательно, 


$\underset{i\in I}{\oplus}{}A_i\in\frak A_C$. 


 


г) Пусть $\{A_i\}_{i\in I}$ --- некоторое семейство групп из класса 


$\frak A_C$, где $I$ --- 


множество неизмеримой мощности, а $C$ --- узкая группа. Тогда, учитывая 


следствие 7 из \cite{2}, имеем 


$\Hom\Big(\prod\limits_{i\in I}A_i,C\Big)\cong 


\underset{i\in I}{\oplus}{}\Hom(A_i,C)=0$. Значит, 


$\prod\limits_{i\in I}A_i\in\frak A_C$. 


 


д) Пусть $\{A_i\;(i\in I);\;\pi^j_i\}$ 


--- прямой спектр абелевых групп и гомоморфизмов (\cite{3}, 


с.\,68), где $A_i\in\frak A_C$. Если $B$ --- подгруппа, порожденная 


всеми элементами вида $a_i-\pi^j_i a_i$ ($i\le j$), то 


$\lim\limits_{\to}{}_I A_i=\underset{i\in I}{\oplus}{}A_i/B$ 


и в силу в) и 


a)\; $\lim\limits_{\to}{}_I A_i\in\frak A_C$. 


 


e) Пусть $A\in\frak A_C$ и $B$ --- произвольная абелева группа. Имеем 


$\Hom(B\otimes A,C)\cong\linebreak \Hom(B,\Hom(A,C))=\Hom(B,0)=0$. 


\end{pf} 


 


\begin{cor} 


Пусть $A$ и $C$ --- абелевы группы, $B$ --- подгруппа группы $A$ такая, 


что $\Hom(B,C)=0$.\; $\Hom(A,C)=0$ тогда и только тогда, когда 


$\Hom(A/B,C)=0$. 


\end{cor} 


 


\begin{pf} 


Необходимость вытекает из п.\,a) предложения 1, а достаточность --- из 


п.\,б) этого предложения. 


\end{pf} 


 


Если $A$ --- абелева группа, то через $T(A)$ обозначим периодическую 


часть группы $A$, а через $T_p(A)$ ($p$ --- простое число) 


$p$-компоненту этой периодической части. Если $A$ --- периодическая 


группа, то ее $p$-компоненту будем обозначать также через $A_p$. Пусть 


$P(A)$ --- множество всех тех простых чисел, для которых $T_p(A)\ne0$. 


Рассмотрим группы гомоморфизмов абелевых групп в периодические группы. 


 


\begin{thm} 


Пусть $A$ и $C$ --- абелевы группы, причем $C$ --- периодическая 


группа.\linebreak $\Hom(A,C)=0$ тогда и только тогда, когда выполняются 


следующие условия$:$


\begin{itemize} 


\item[1)] если $C$ --- нередуцированная группа, то $A$ --- периодическая 


группа$;$ 


\item[2)] для всякого $p\in P(A)\cap P(C)$\; $T_p(A)$ --- делимая, 


$C_p$ --- редуцированная группы$;$ 


\item[3)] для всякого $p\in P(C)$ факторгруппа $A/T(A)$\; $p$-делима. 


\end{itemize} 


\end{thm} 


 


\begin{pf} 


{\bf Необходимость.} Пусть $\Hom(A,C)=0$, где $C$ --- периодическая 


группа. Предположим, что $C$ --- нередуцированная группа, и пусть группа 


$A$ содержит элементы 


бесконечного порядка. Если $V$ --- делимая часть группы $C$ и $a\in A$, 


$o(a)=\infty$, то $\Hom(\la a\ra,V)\ne0$. В силу инъективности группы 


$V$ любой гомоморфизм 


подгруппы $\la a\ra$ группы 


$A$ в группу $V$ можно продолжить до гомоморфизма $A$ в $V$. 


Следовательно, 


$\Hom(A,V)\ne0$, и поэтому $\Hom(A,C)\ne0$. Противоречие. Значит, 


условие~1) выполняется. 


 


Пусть $q\in P(A)\cap P(C)$. Если $C_q$ --- нередуцированная группа, то, 


рассматривая в группе 


$A_q$ произвольную ненулевую циклическую подгруппу $\la a\ra$, получим 


$\Hom(\la a\ra,V_q)\ne0$, где $V_q$ --- делимая часть группы $C_q$. 


Отсюда имеем $\Hom(A,V_q)\ne0$, и, значит, 


$\Hom(A,C)\ne0$. Итак, $C_q$ --- редуцированная группа. Предположим, 


что $T_q(A)$ имеет ненулевую 


редуцированную часть. Имеем $T_q(A)=R_1\oplus V_1$, где $R_1$, $V_1$ 


--- редуцированная и делимая 


части группы $T_q(A)$ соответственно. Группы $R_1$ и $C_q$ обладают 


циклическими 


прямыми слагаемыми $B_1$ и $B_2$ соответственно, порядки которых суть 


степени простого числа $q$. Так как $B_1$ --- сервантная циклическая 


$q$-подгруппа группы $A$, то $B_1$ --- прямое слагаемое группы $A$. 


Имеем 


\begin{multline*} 


\Hom(A,C_q)=\\ =\Hom(B_1\oplus B'_1,\;B_2\oplus B'_2)\cong 


\Hom(B_1,B_2)\oplus\Hom(B_1,B'_2)\oplus 


\Hom(B'_1,B_2)\oplus\Hom(B'_1,B'_2)\ne 0 


\end{multline*}  


т.\,к.~$\Hom(B_1,B_2)\ne0$. Отсюда $\Hom(A,C)\ne0$. Значит, $T_q(A)$ 


--- делимая группа, и условие~2) выполняется. 


 


Итак, мы показали, что если $\Hom(A,C)=0$, 


то условия 1) и 2) выполняются. При этом 


имеем $\Hom(T(A),C)=0$. Действительно, 


$\Hom(T(A),C)\cong\prod\limits_p\Hom(T_p(A),C_p)=0$, т.\,к.~группа 


гомоморфизмов делимой группы 


$T_p(A)$ в редуцированную группу $C_p$ равна нулю. Тогда, применяя 


следствие 


1, получим $\Hom(A/T(A),C)=0$. Предположим, что для некоторого 


$q\in P(C)$ факторгруппа $A/T(A)$ не 


является $q$-делимой. Пусть $\ov a\in A/T(A)$, $\ov a$ имеет нулевую 


$q$-высоту, и пусть $c$ 


--- произвольный ненулевой элемент группы $C_q$. Рассмотрим гомоморфизм 


$\varphi:\la\ov a\ra_*\to\la c\ra$ 


такой, что $\varphi\ov a=c$. Так как $\la c\ra$ --- алгебраически 


компактная группа, то $\varphi$ 


продолжается до гомоморфизма $\ov\varphi:A/T(A)\to\la c\ra$.\; 


$\ov\varphi\ne0$, и поэтому $\Hom(A/T(A),\la c\ra\ne0$, а значит, и 


$\Hom(A/T(A),C)\ne0$. Итак, условие~3) выполняется. 


 


{\bf Достаточность.} Пусть выполняются условия 1)--3). В силу 


условия~2)\; $\Hom(T(A),C)=0$. Покажем, что и $\Hom(A/T(A),C)=0$; тогда 


по следствию~1\; $\Hom(A,C)=0$. Если $C$ 


--- нередуцированная группа, то по условию~1)\; $A$ --- периодическая 


группа, и поэтому $\Hom(A/T(A),C)=\Hom(0,C)=0$. Если же $C$ --- 


редуцированная группа, то каждая ее подгруппа $C_p$ ($p\in P(C)$\,) 


является редуцированной. Предположим, что 


$\Hom(A/T(A),C_p)\ne0$, и пусть $\varphi\in\Hom(A/T(A),C_p)$, 


$\varphi\ne0$. В силу условия~3)\; $A/T(A)$ 


--- $p$-делимая группа, тогда и $\im\varphi$ --- 


$p$-делимая группа, и, значит, в $C_p$ есть ненулевые делимые подгруппы, 


чего быть не может. Итак $\Hom(A/T(A),C_p)=0$ для каждого $p\in P(C)$. 


Значит, 


$\Hom\Big(A/T(A),\prod\limits_{p\in P(C)}C_p\Big)\cong 


\prod\limits_{p\in P(C)}\Hom(A/T(A),C_p)=0$. Так как 


$\underset{p\in P(C)}{\oplus}{}C_p$ --- 


подгруппа группы $\prod\limits_{p\in P(C)}C_p$, то 


$\Hom(A/T(A),C)=\Hom(A/T(A),\underset{p\in P(C)}{\oplus}{}C_p)=0$. 


\end{pf} 


 


\begin{cor} 


Пусть $A$ --- абелева группа без кручения, $C$ --- периодическая абелева 


группа.\; $\Hom(A,C)=0$ тогда и только тогда, когда $C$ --- 


редуцированная группа и для всякого $p\in P(C)$ группа $A$\; $p$-делима. 


\end{cor} 


 


\begin{cor} 


Пусть $A$ --- смешанная абелева группа, $C$ --- периодическая абелева 


группа.\; $\Hom(A,C)=0$ тогда и только тогда, когда выполняются 


следующие условия: 


\begin{itemize} 


\item[1)] $C$ --- редуцированная группа; 


\item[2)] для всякого $p\in P(A)\cap P(C)$\; $T_p(A)$ 


--- делимая группа; 


\item[3)] для всякого $p\in P(C)$ факторгруппа $A/T(A)$\; $p$-делима.


\end{itemize} 


\end{cor} 


 


Учитывая изоморфизм 


$\Hom(A,C)\cong\prod\limits_p\Hom(A_p,T_p(C)$, где $A$ --- 


периодическая, а $C$ --- произвольная абелевы группы, получаем 


 


\begin{cor} 


Пусть $A$ и $C$ --- абелевы группы, причем $A$ --- периодическая 


группа.\linebreak


$\Hom(A,C)=0$ тогда и только тогда, когда для всякого 


$p\in P(A)\cap P(C)$\; $A_p$ --- делимая 


группа, $T_p(C)$ --- редуцированная группа. 


\end{cor} 


 


Рассмотрим теперь группы гомоморфизмов абелевых групп в абелевы группы 


без кручения ранга~1. 


 


\begin{thm} 


Пусть $C$ --- абелева группа без кручения ранга $1$ типа $t$, $A$ --- 


абелева группа, тип каждого ненулевого элемента факторгруппы $A/T(A)$ 


которой 


не меньше $t$.\; $\Hom(A,C)=0$ тогда и только тогда, когда группа 


$A/T(A)$ не содержит прямого слагаемого, изоморфного $C$. 


\end{thm} 


 


\begin{pf} 


{\bf Необходимость.} Пусть $\Hom(A,C)=0$ и предположим противное, что 


группа 


$A/T(A)$ содержит прямое слагаемое $A_1/T(A)$, изоморфное группе $C$. 


Имеем $A/T(A)=A_1/T(A)\oplus A_2/T(A)$. Тогда 


$$ 


\Hom(A/T(A),C)\cong\Hom(A_1/T(A)\oplus A_2/T(A),C)\cong 


\Hom(A_1/T(A),C)\oplus\Hom(A_2/T(A),C)\ne 0. 


$$ 


В силу следствия 1\; $\Hom(A,C)\ne0$. Противоречие. 


 


{\bf Достаточность.} Обозначим группу $A/T(A)$ через $\ov A$. Пусть 


группа $\ov A$ не 


содержит прямого слагаемого, изоморфного $C$, и предположим, что 


$\Hom(A,C)\ne0$. В 


силу следствия 1\; $\Hom(\ov A,C)\ne0$. Тогда существует 


$f\in\Hom(\ov A,C)$, $f\ne0$. Найдется такой элемент 


$\ov a\in\ov A$, что $f(\ov a)\ne0$. Так как при гомоморфизме типы 


элементов не уменьшаются, 


то, учитывая условие на типы элементов группы $\ov A$, получаем 


$t(\ov a)=t(f(\ov a))=t$. 


Характеристика элемента $f(\ov a)$ в группе $\im f$ не меньше 


характеристики элемента 


$\ov a$ и поэтому $t(\im f)=t$. Итак, $\im f\cong C$. Имеем 


$\ov A\Ker f\cong C$.\; $\Ker f$ --- 


сервантная подгруппа 


группы $\ov A$, и все элементы множества $\ov A\setminus\Ker f$ 


имеют тип $t$. Тогда (\cite{4}, предложение 86.5, с.\,136) $\Ker f$ --- 


прямое слагаемое группы $\ov A$, дополнительное же прямое слагаемое 


изоморфно $C$. Получили противоречие. 


\end{pf} 


 


Заметим, что в работе \cite{1} абелевы группы $A$ со свойством 


$\Hom(A,Z)=0$ названы коузкими. 


 


Напомним, что абелева группа $A$ называется $p$-локальной ($p$ --- 


простое число), если $nA=A$ для всех целых чисел, взаимно простых 


с $p$. 


 


\begin{cor} 


Пусть $A$ --- $p$-локальная абелева группа, $C$ --- редуцированная 


абелева группа без кручения ранга 1.\; $\Hom(A,C)=0$ тогда и только 


тогда, когда группа $A/T(A)$ 


не содержит прямого слагаемого, изоморфного $C$. 


\end{cor} 


 


\begin{pf} 


Тип группы $A/T(A)$ не меньше типа любой редуцированной абелевой группы 


без кручения ранга 1. Остается применить теорему 2. 


\end{pf} 


 


\begin{cor} 


Пусть $A$ --- абелева группа, решетка вполне характеристических подгрупп 


которой является цепью, $C$ --- редуцированная абелева группа без 


кручения ранга 1.\linebreak $\Hom(A,C)=0$ тогда и только тогда, когда группа 


$A/T(A)$ не содержит прямого слагаемого, изоморфного $C$. 


\end{cor} 


 


\begin{pf} 


Всякая абелева группа, решетка вполне характеристических подгрупп 


которой образует цепь, является $p$-локальной группой \cite{5}. Теперь 


применим следствие 5. 


\end{pf} 


 


Пусть $t$ --- некоторый тип (класс эквивалентных характеристик). 


Обозначим через $\frak M_t$ класс всех таких абелевых групп, которые 


не содержат прямого слагаемого без кручения ранга 1 типа $t$, а через 


$\frak N_t$ --- класс 


всех таких абелевых групп без кручения, у которых тип каждого ненулевого 


элемента не меньше $t$. 


 


\begin{thm} 


Пусть $A$, $A'$, $A_i$ $(i\in I)$ --- абелевы группы такие, что 


$A/T(A)$, $A'/T(A')$, $A_i/T(A_i)$ $(i\in I)$ принадлежат классу 


$\frak M_t\cap\frak N_t$. 


 


$1.$ Следующие группы принадлежат классу $\frak M_t{:}$ {\em a)}~любая 


факторгруппа группы


$A$ $($в частности, сама группа $A);$ {\em б)}~расширение группы $A$ с 


помощью $A';$ {\em в)}~прямая сумма групп $A_i$ $(i\in I);$ 


{\em г)}~прямое 


произведение групп $A_i$ $(i\in I)$, если множество $I$ неизмеримо$;$ 


{\em д)}~предел прямого спектра, построенного на группах $A_i$ 


$(i\in I);$ {\em e)}~тензорное произведение $D\otimes A$, где $D$ --- 


произвольная абелева группа. 


 


$2.$ Каждая из групп, получающаяся в пп.\,{\em a)--e)}, обладает тем 


свойством, что ее группа гомоморфизмов в любую однородную сепарабельную 


группу без кручения типа $t$ $($в частности, в абелеву группу без 


кручения ранга $1$ типа $t)$ равна нулю. 


\end{thm} 


 


\begin{pf} 


Пусть $C$ --- группа без кручения ранга 1 типа $t$. Так как $A/T(A)$, 


$A'/T(A')$, $A_i/T(A_i)$ ($i\in I$) --- группы из класса 


$\frak M_t\cap\frak N_t$, то по 


теореме 2 каждая из групп гомоморфизмов $\Hom(A,C)$, $\Hom(A',C)$, 


$\Hom(A_i,C)$ равна нулю. Тогда по предложению 1 группа 


гомоморфизмов каждой из групп, получающихся в пп.\,a)--e), в группу $C$ 


равна нулю. Значит, каждая из групп, получающаяся в пп.\,a)--e), 


принадлежит классу $\frak M_t$. 


 


Пусть теперь $B$ --- одна из групп, получающаяся в пп.\,a)--e), а $C'$ 


--- однородная сепарабельная группа типа $t$. Группа $C'$ изоморфна 


некоторой 


сервантной подгруппе группы $\Big(\prod\limits_i R_i\Big)(t)$, 


где группы $R_i$ имеют ранг 1 и 


тип $t$ (\cite{4}, предложение 87.4, с.\,142). Рассмотрим 


$\Hom\Big(B,\prod\limits_i R_i\Big)$. Имеем 


$\Hom\Big(B,\prod\limits_i R_i\Big)\cong 


\prod\limits_i\Hom(B,R_i)=0$, т.\,к.~для всякой группы $R_i$\; 


$\Hom(B,R_i)=0$. Значит, $\Hom(B,C')=0$. 


\end{pf} 


 


Обозначим через $\frak L_t$ класс всех однородных абелевых групп без 


кручения типа $t$, которые не содержат прямых слагаемых ранга 1. Класс 


$\frak L_t$ является 


подклассом класса $\frak M_t\cap\frak N_t$. 


 


Из теоремы 3 вытекает такое следствие. 


 


\begin{cor} 


1. Следующие теоретико-групповые конструкции групп из $\frak L_t$ не 


содержат 


прямых слагаемых ранга 1 типа $t$: a)~факторгруппы; б)~расширения; 


в)~прямые суммы; г)~прямые произведения с неизмеримым множеством 


компонент; д)~прямые пределы; e)~тензорные произведения на произвольную 


абелеву группу. 


 


2. Каждая из групп, получающаяся в пп.\,a)--e) (в том числе и любая 


группа из $\frak L_t$), обладает тем свойством, что ее группа 


гомоморфизмов в 


любую однородную сепарабельную группу типа~$t$ (в частности, в абелеву 


группу без кручения ранга 1 типа $t$) равна нулю. 


\end{cor} 


 


\begin{cor} 


Пусть $A$ и $C$ --- однородные абелевы группы без кручения одного и того 


же типа, причем $C$ --- сепарабельная группа.\; $\Hom(A,C)=0$ тогда и 


только тогда, когда $A$ не содержит прямого слагаемого ранга 1. 


\end{cor} 


 


\begin{cor} 


Группа гомоморфизмов любой алгебраически компактной группы также, как и 


любой копериодической группы, в любую редуцированную однородную 


сепарабельную группу без кручения (в частности, в любую редуцированную 


группу без кручения ранга 1) равна нулю. 


\end{cor} 


 


\begin{pf} 


Пусть $A$ --- алгебраически компактная группа. Не умаляя общности, можно 


считать, что $A$ --- редуцированная группа. Имеем 


$A=\prod\limits_p A_p$, где каждая группа 


$A_p$ полна в своей $p$-адической топологии (произведение берется по 


различным простым числам $p$).\; $A_p/T(A_p)$ --- алгебраически 


компактная группа 


(\cite{3}, с.\,90), и т.\,к.~не существует счетных редуцированных 


алгебраически компактных групп без кручения (\cite{3}, следствие 40.4, 


с.\,199), то $A_p/T(A_p)$ не содержит редуцированного прямого слагаемого 


без кручения ранга 1. Тогда в силу теоремы 3 для любой редуцированной 


однородной сепарабельной группы без кручения $C$\; $\Hom(A,C)=0$. 


 


Каждая копериодическая группа есть эпиморфный образ алгебраически 


компактной группы, и поэтому, применяя предложение 1 (п.\,a)\,), 


получим, что группа гомоморфизмов любой копериодической группы в любую 


редуцированную однородную сепарабельную группу без кручения равна нулю. 


\end{pf} 


 


\begin{thebibliography}{9} 


 


\bibitem{1} 


Dimitri\v c R. {\em On coslender groups} // Glasnik Matem. -- 1986. -- 


V.\,21. -- \No\,2. -- P.\,327--329. 


\bibitem{2} 


Рычков С.В. {\em О прямых произведениях абелевых групп} // Матем. сб. -- 


1982. -- Т.\,117. -- \No\,2. -- С.\,266--278. 


\bibitem{3} 


Фукс Л. {\em Бесконечные абелевы группы}. Т.\,1. -- М.: Мир, 1974. -- 


335 с. 


\bibitem{4} 


Фукс Л. {\em Бесконечные абелевы группы}. Т.\,2. -- М.: Мир, 1977. -- 


416 с. 


\bibitem{5} 


Hausen J. {\em Abelian groups which are uniserial as modules over their 


endomorphism rings} // Lect. Notes Math. -- 1983. -- V.\,1006. -- 


P.\,204-208. 


 


\end{thebibliography} 


 


\vskip 2cm 


 


\postup{\it Томский государственный}{$\qquad$\it Поступили} 


\postup{\it университет}{{\it первый вариант} 02.02.1996} 


\postup{}{{\it окончательный вариант} 26.03.1997} 


 


\label{end} 


\end{document} 


